
    The input consists of:
    \begin{itemize}
        \item One line with either ``\texttt{Encode}'' or ``\texttt{Decode}'', depending on the action you are to perform.
        \item One line with an integer $n$ ($1 \leq n \leq 10^5$), the size of the data packet.
        \item One line with a string $s$ of $n$ characters, the data packet.
    \end{itemize}

    In the encoding pass, this string consists of the symbols `\texttt{0}' and `\texttt{1}'.

    In the decoding pass, this string will be a cyclic rotation of your ternary string. In other words, this string is the string you printed in the encoding pass, potentially modified to reflect how this data packet might have been received on Earth, as described in the statement.
